\documentclass[12pt,a4paper]{article}

\usepackage[utf8]{inputenc}
\usepackage[T1]{fontenc}
\usepackage[brazilian]{babel}
\usepackage[margin=2.5cm]{geometry}
\usepackage{setspace}
\usepackage{indentfirst}
\usepackage{hyperref}
\usepackage{url}
\usepackage{enumitem}

\onehalfspacing

\hypersetup{
  colorlinks=true,
  linkcolor=black,
  citecolor=black,
  urlcolor=blue,
  pdftitle={Proposta de Mestrado em Sonologia},
}

\title{\textbf{wasmwasm: linguagens de programação musical no navegador \\
como mediação do pensamento composicional} \\[0.3cm]
\large Proposta de Projeto de Pesquisa para Seleção de Mestrado \\
Programa de Pós-Graduação em Música da UFMG \\
Área de Concentração: Sonologia \\
Linha de Pesquisa: Música e Tecnologia}

\author{[Nome do candidato] \\ \small e-mail: [a completar] \\ \small Possível orientação: Prof. Dr. Sérgio Freire (ou docente da linha, a confirmar)}

\date{2026}

\begin{document}

\maketitle

\begin{abstract}
\noindent Esta proposta parte do desenvolvimento do \emph{wasmwasm}, um ambiente de síntese sonora e composição algorítmica que roda inteiramente no navegador, compilando duas linguagens de domínio específico (DSLs) - uma linguagem de \emph{patch}, para definição de instrumentos/vozes sonoras, e uma linguagem de \emph{score}, para sequenciamento e organização temporal - para WebAssembly, com execução em tempo real via Web Audio API. A pesquisa proposta investiga como escolhas de design dessas linguagens (tipagem, separação entre instrumento e partitura, primitivas de controle temporal) moldam, habilitam e constrangem o pensamento composicional de quem as utiliza. A metodologia combina pesquisa-ação/pesquisa por design - o desenvolvimento contínuo do sistema como parte do processo investigativo - com um estudo de caso qualitativo envolvendo compositores externos em oficinas de uso da ferramenta. O trabalho articula-se com a tradição de linguagens de programação musical (Csound, SuperCollider, ChucK, TidalCycles) e com os estudos de \emph{live coding} e de linguagem como mediação técnica do fazer musical, contribuindo para a sonologia com uma reflexão situada sobre autoria, ferramenta e pensamento sonoro no contexto de tecnologias nativas de navegador.

\vspace{0.3cm}
\noindent \textbf{Palavras-chave:} sonologia; linguagens de programação musical; WebAssembly; composição algorítmica; live coding; síntese sonora.
\end{abstract}

\section{Introdução e Justificativa}

O ato de compor com código é, ao mesmo tempo, um ato de programar e um ato de pensar musicalmente através de uma linguagem formal. Desde o Music~N de Max Mathews até ambientes contemporâneos como Csound, SuperCollider, ChucK e TidalCycles, a história da música eletroacústica e da sonologia é indissociável da história das linguagens criadas para descrever som e tempo computacionalmente. Cada uma dessas linguagens carrega, em suas escolhas sintáticas e semânticas, um modelo implícito de como o som pode ser pensado, organizado e manipulado - e, por extensão, molda o próprio pensamento composicional de quem a utiliza.

O projeto \emph{wasmwasm}\footnote{Repositório do projeto: \url{https://github.com/guilibre/wasmwasm} (a confirmar/ajustar conforme publicação).}, desenvolvido pelo candidato, insere-se nesse território, com uma particularidade tecnológica relevante: trata-se de um ambiente de síntese e composição que roda inteiramente no navegador, sem necessidade de instalação, compilando código-fonte para WebAssembly (WASM) e executando síntese em tempo real através da Web Audio API (mais especificamente, via \texttt{AudioWorkletProcessor}). O sistema é composto por dois compiladores, escritos em C++ e compilados via Emscripten:

\begin{itemize}
  \item uma \textbf{linguagem de patch}, para definição de instrumentos e vozes sonoras, com inferência de tipos ao estilo Hindley-Milner e geração de código intermediário otimizado - incluindo \emph{unrolling} automático de laços para integradores numéricos de passo fixo, usados tipicamente em filtros e envelopes;
  \item uma \textbf{linguagem de score}, para sequenciamento e organização temporal, cujo código-fonte é resolvido em um grafo de nós (\texttt{state}, \texttt{fork}, \texttt{join}, \texttt{transform}, \texttt{legato}, \texttt{repeat}, \texttt{reverse}, \texttt{skip}, entre outros), interpretado em tempo real por um módulo \emph{conductor} que percorre o grafo e aciona instâncias de instrumentos WASM dentro do \texttt{AudioWorklet}.
\end{itemize}

O frontend, em React/TypeScript, oferece editor de código com suporte a \emph{Language Server Protocol} (LSP), osciloscópio, espectrograma, visualização em grafo da partitura e piano roll - reunindo, em um único ambiente acessível via navegador, ferramentas de escrita, execução e visualização sonora.

Essa arquitetura - duas linguagens decopladas, uma para o timbre/instrumento e outra para a organização temporal/partitura, ambas compiladas e executadas localmente no navegador do usuário - constitui um caso de estudo pouco explorado na literatura de sonologia e de linguagens de programação musical: a maior parte dos ambientes canônicos (Csound, SuperCollider, ChucK) foi concebida para execução nativa, com modelos de instalação, texto único ou paradigmas de live coding centrados em uma linguagem só. A proposta de investigar como a separação patch/score e as escolhas de tipagem e de primitivas de grafo do \emph{wasmwasm} afetam o pensamento composicional de seus usuários contribui tanto para a prática de desenvolvimento de ferramentas musicais quanto para a reflexão teórica sobre linguagem como mediadora da criação sonora - núcleo de interesse da sonologia enquanto campo que articula tecnologia, escuta e composição.

\section{Problema e Pergunta de Pesquisa}

A pergunta central que orienta esta proposta é:

\begin{quote}
\emph{De que modo as escolhas de design de uma linguagem de programação musical - nomeadamente, a separação entre linguagem de instrumento (patch) e linguagem de partitura (score), o sistema de tipos e as primitivas de controle temporal do grafo de score - moldam, habilitam e constrangem o pensamento composicional de compositores que a utilizam?}
\end{quote}

Desdobram-se dessa pergunta central três questões secundárias:

\begin{enumerate}
  \item Como a separação arquitetural entre a linguagem de instrumento, voltada a timbre e síntese, e a linguagem de partitura, voltada a tempo e sequenciamento, refletida nas duas DSLs do \emph{wasmwasm}, se compara a modelos de outras linguagens musicais, e que efeitos essa separação tem sobre o processo composicional?
  \item Que primitivas do grafo de score (fork, join, transform, legato, repeat, reverse, skip) emergem como recursos expressivos genuinamente novos ou como reformulações de conceitos musicais já conhecidos (ex.: canon, retrogradação, transformação paramétrica)?
  \item Em que medida a execução no navegador, sem instalação, e a natureza acessível/compartilhável do ambiente afetam quem consegue experimentar composição algorítmica e síntese sonora, e como isso se relaciona à apropriação criativa da linguagem?
\end{enumerate}

\section{Objetivos}

\subsection{Objetivo Geral}

Investigar, por meio do desenvolvimento contínuo do sistema \emph{wasmwasm} e de sua utilização por compositores, como as escolhas de design de suas linguagens de patch e de score moldam o pensamento e o processo composicional de quem as utiliza, contribuindo para a reflexão sonológica sobre linguagem como mediação técnica da criação musical.

\subsection{Objetivos Específicos}

\begin{itemize}
  \item Consolidar e ampliar o desenvolvimento técnico do \emph{wasmwasm} durante o mestrado, documentando decisões de design das linguagens de patch e de score à luz da literatura de linguagens de programação musical;
  \item Situar teoricamente o \emph{wasmwasm} face a ambientes correlatos (Csound, SuperCollider, ChucK, TidalCycles, Faust, Sporth, entre outros), identificando continuidades e rupturas;
  \item Conduzir um estudo de caso qualitativo - oficina(s) com compositores externos ao projeto usando as DSLs do \emph{wasmwasm} - para observar e analisar como a linguagem afeta suas decisões e seu pensamento composicional;
  \item Produzir um corpus de exemplos e/ou obras (próprias e/ou dos participantes da oficina) que evidenciem usos e apropriações da ferramenta;
  \item Articular, a partir dos resultados empíricos e do processo de desenvolvimento, uma reflexão teórica sobre a relação entre design de linguagem, materialidade computacional e pensamento sonoro/composicional.
\end{itemize}

\section{Referencial Teórico}

A fundamentação teórica desta proposta articula-se em três eixos principais, a serem aprofundados ao longo do mestrado:

\textbf{(1) Linguagens de programação musical e sonologia.} A literatura sobre ambientes de síntese e composição algorítmica - de Curtis Roads (\emph{Microsound}; \emph{Computer Music Tutorial}) às análises comparativas de linguagens como Csound, SuperCollider e ChucK - oferece o pano de fundo técnico-histórico para situar as escolhas do \emph{wasmwasm}.

\textbf{(2) Live coding e linguagem como escrita sonora.} A obra de Thor Magnusson, em especial \emph{Sonic Writing: Technologies for Composing, Performing and Notating Music} (2019), é central para pensar como sistemas de notação e linguagens de programação musical funcionam como tecnologias de escrita do som, moldando cognição e prática musical. Os trabalhos de Nick Collins e da comunidade de \emph{live coding} (incluindo o TOPLAP e publicações associadas ao \emph{International Conference on Live Coding}) fornecem referências sobre a relação entre código, tempo real e pensamento musical emergente.

\textbf{(3) Interfaces, mediação técnica e design de instrumentos musicais digitais.} O trabalho de Andrew McPherson (e colaborações com Magnusson) sobre o design de \emph{Digital Musical Instruments} e sobre a relação entre affordances técnicas e prática musical fundamenta a análise de como decisões de engenharia (tipagem, arquitetura patch/score, execução em WASM/navegador) se traduzem em possibilidades e limitações expressivas.

Complementarmente, a proposta dialoga com a filosofia da técnica aplicada à música (mediação técnica, não-neutralidade da ferramenta) e com estudos sobre acessibilidade e democratização de ferramentas criativas digitais, na medida em que a execução no navegador, sem instalação, constitui uma escolha técnica com implicações sobre quem pode compor e experimentar com síntese sonora.

\section{Metodologia}

A pesquisa proposta é de natureza prático-teórica, combinando duas abordagens metodológicas complementares:

\subsection{Pesquisa-ação / pesquisa por design}

O desenvolvimento contínuo do \emph{wasmwasm} ao longo do mestrado constitui parte central do método: decisões de design das linguagens de patch e de score serão tomadas, documentadas e revisadas como parte do processo investigativo, seguindo uma lógica de pesquisa por design (\emph{research through design}), na qual o artefato construído é também instrumento e resultado de conhecimento.

\subsection{Estudo de caso qualitativo}

Serão conduzidas uma ou mais oficinas com compositores externos ao projeto, convidados a utilizar as DSLs do \emph{wasmwasm} para compor peças curtas ou estudos sonoros. Os dados a coletar incluem:

\begin{itemize}
  \item registros do código produzido pelos participantes (análise de padrões de uso, dificuldades, apropriações não previstas);
  \item entrevistas semiestruturadas antes/depois da oficina, sobre processo composicional e percepção da linguagem;
  \item observação participante durante as sessões de uso.
\end{itemize}

A análise dos dados qualitativos seguirá procedimentos de análise temática, articulando achados empíricos com o referencial teórico apresentado na Seção 4.

\section{Cronograma}

Cronograma estimado para 24 meses, estrutura típica de mestrado no PPGMUS-UFMG:

\begin{center}
\begin{tabular}{|p{7cm}|c|c|c|c|c|c|c|c|}
\hline
\textbf{Atividade} & \multicolumn{8}{c|}{\textbf{Trimestre}} \\
\hline
 & 1 & 2 & 3 & 4 & 5 & 6 & 7 & 8 \\
\hline
Revisão bibliográfica e créditos & X & X & X & & & & & \\
\hline
Desenvolvimento do \emph{wasmwasm} (contínuo) & X & X & X & X & X & X & & \\
\hline
Desenho e preparação da(s) oficina(s) & & X & X & & & & & \\
\hline
Realização da(s) oficina(s) / coleta de dados & & & X & X & & & & \\
\hline
Análise dos dados qualitativos & & & & X & X & & & \\
\hline
Redação da dissertação & & & & X & X & X & X & \\
\hline
Qualificação & & & & & X & & & \\
\hline
Revisão final e defesa & & & & & & & X & X \\
\hline
\end{tabular}
\end{center}

\section{Resultados Esperados}

\begin{itemize}
  \item Versão consolidada e documentada do sistema \emph{wasmwasm}, com registro das decisões de design tomadas durante o mestrado;
  \item Corpus de composições e/ou estudos sonoros produzidos com a ferramenta, pelo candidato e/ou participantes da oficina;
  \item Dissertação de mestrado articulando desenvolvimento técnico, estudo de caso empírico e reflexão teórica sobre linguagem de programação musical como mediação do pensamento composicional;
  \item Publicações em eventos/periódicos da área (ex.: Simpósio Brasileiro de Computação Musical, NIME, International Conference on Live Coding), conforme oportunidade e maturidade dos resultados.
\end{itemize}

\section*{Referências (fundação inicial)}
\addcontentsline{toc}{section}{Referências}

\begin{itemize}[leftmargin=1.5cm]
  \item COLLINS, Nick; MCLEAN, Alex; ROHRHUBER, Julian; WARD, Adrian. Live coding in laptop performance. \emph{Organised Sound}, v. 8, n. 3, p. 321–330, 2003.
  \item MAGNUSSON, Thor. \emph{Sonic Writing: Technologies for Composing, Performing and Notating Music}. New York: Bloomsbury Academic, 2019.
  \item MCPHERSON, Andrew; ZAPPI, Victor. An environment for submillisecond-latency audio and sensor processing on BeagleBone Black. In: \emph{Audio Engineering Society Convention}, 2015.
  \item ROADS, Curtis. \emph{The Computer Music Tutorial}. Cambridge, MA: MIT Press, 1996.
  \item ROADS, Curtis. \emph{Microsound}. Cambridge, MA: MIT Press, 2001.
  \item WANG, Ge. The ChucK Audio Programming Language: A Strongly-timed and On-the-fly Environ/mentality. Tese (Doutorado) - Princeton University, 2008.
  \item MCCARTNEY, James. Rethinking the computer music language: SuperCollider. \emph{Computer Music Journal}, v. 26, n. 4, p. 61–68, 2002.
  \item VUOSKOSKI, Jonna K.; MAGNUSSON, Thor et al. (referências complementares sobre cognição musical e linguagem, a levantar).
  \item [Referências adicionais sobre WebAssembly, Web Audio API e síntese sonora em navegador, a levantar durante a revisão bibliográfica.]
\end{itemize}

\end{document}
